% Part: normal-modal-logic
% Chapter: filtrations
% Section: filtrations-def

\documentclass[../../../include/open-logic-section]{subfiles}

\begin{document}

\olfileid{nml}{fil}{fil}

\olsection{వడపోతలు}

$\Gamma$ ద్వారా $\mModel{M}$కు వడపోత ఒకటే అన్నట్లు దాన్ని
నిర్వచించకుండా, ఒక నమూనా~$\mModel{M^*}$ను ఎప్పుడు
$\mModel{M}$కు వడపోత అవుతుందో నిర్వచిస్తాం. అదే నమూనాకు అదే సమితి ద్వారా చేసే వడపోతలన్నింటికీ
లోకాల సమితి~$W^*$, మూల్యనిర్ణయం~$V^*$ ఒకటే. కానీ వేర్వేరు వడపోతల్లో
ప్రాప్యత సంబంధం~$R^*$ మారవచ్చు.
వడపోతగా పరిగణించాలంటే $R^*$ కొన్ని షరతులను
నెరవేర్చాలి. ప్రధాన ఫలితం, అంటే $!A \in \Gamma$
అయినప్పుడు $\mSat{M}{!A}[w]$ నెరవేరేది
$\mSat{M^*}{!A}[{[w]}]$ నెరవేరినప్పుడు, అప్పుడే మాత్రమే అని నిరూపించడానికి
అవసరమైన షరతులే అవి.

\begin{defn}\ollabel{defn:filtration}
  $\Gamma$ ఉపసూత్రాల పరంగా సంవృతమై ఉండనివ్వండి;
  $\mModel{M} = \tuple{W, R, V}$ అనుకుందాం.
  $\Gamma$ ద్వారా $\mModel{M}$కు చేసే \emph{వడపోత}
  అంటే, ఈ విధమైన ఏ నమూనా $\mModel{M^*} = \tuple{W^*,R^*,V^*}$ అయినా:
  \begin{enumerate}
  \item $W^* = \Setabs{[w]}{w \in W}$;
  \item \ollabel{defn:filtration-R}%
    ఏ $u,v \in W$కైనా:
    \begin{enumerate}
    \item \ollabel{defn:filtration-R1}%
      $Ruv$ అయితే $R^*[u][v]$;
    \tagitem{prvBox}{\ollabel{defn:filtration-R2}%
      $R^*[u][v]$ అయితే, ఏ $\Box !A \in \Gamma$కైనా,
      $\mSat{M}{\Box!A}[u]$ అయినప్పుడు
      $\mSat{M}{!A}[v]$}{};
    \tagitem{prvDiamond}{\ollabel{defn:filtration-R3}%
      $R^*[u][v]$ అయితే, ఏ $\Diamond !A \in \Gamma$కైనా,
      $\mSat{M}{!A}[v]$ అయినప్పుడు
      $\mSat{M}{\Diamond!A}[u]$.}{}
    \end{enumerate}
  \item $V^*(p) = \Setabs{[u]}{u \in V(p)}$.
  \end{enumerate}
\end{defn}

$V^*(p)$లో ఏయే వర్గాలు ఉంటాయో చూద్దాం.
$\mModel{M}$లో $p$ సత్యమయ్యే ప్రతి లోకం~$w$కూ
దాని తుల్యతా వర్గం~$[w]$ను తీసుకుంటాం. అలా వచ్చే
వర్గాలే ఈ సమితిలోని మూలకాలు.
ఒకవైపు, $w \in V(p)$ అయితే ఈ నిర్వచనం వల్ల
$[w] \in V^*(p)$. అయితే $[w] \in V^*(p)$ అయితే
$w \in V(p)$ తప్పనిసరి కాదు. $[w] \in V^*(p)$
అయినప్పుడు \emph{ఏదో}~$u \in V(p)$కు $[w] = [u]$
అని మాత్రమే చెప్పగలం. $[w] = [u]$ అంటే $w \equiv u$.
కాబట్టి $[w] \in V^*(p)$ నుంచి, ఏదో $u \in
V(p)$కు $w \equiv u$ అని మాత్రమే తేల్చగలం.

\begin{thm}\ollabel{thm:filtrations}
  $\mModel{M^*}$ అనేది $\Gamma$ ద్వారా
  $\mModel{M}$కు చేసే వడపోత అయితే, ప్రతి
  $!A \in \Gamma$కూ, ప్రతి $w \in W$కూ,
  $\mSat{M}{!A}[w]$ నెరవేరేది
  $\mSat{M^*}{!A}[{[w]}]$ నెరవేరినప్పుడు, అప్పుడే మాత్రమే.
\end{thm}

\begin{proof}
  $\Gamma$ ఉపసూత్రాల పరంగా సంవృతమనే విషయాన్ని
  ఉపయోగించి $!A$ నిర్మాణంపై ఆగమనం చేద్దాం. $!A \in \Gamma$,
  $\Gamma$ \tetoken{ఉపసూత్రాల}{!!{formula}s} పరంగా
  సంవృతం కనుక~$!A$కు చెందిన అన్ని
  \tetoken{ఉపసూత్రాలు}{!!{formula}s} కూడా $\in
  \Gamma$. అందువల్ల ప్రతి ఆగమన దశలో~$!A$కు చెందిన
  \tetoken{ఉపసూత్రాలకు}{!!{formula}s} ఆగమన పరికల్పన
  వర్తిస్తుంది.
  \begin{enumerate}
    \tagitem{prvFalse}{\indcase{!A}{\lfalse}{$\mSat{M}{\indfrm}[w]$
        గానీ $\mSat{M^*}{\indfrm}[{[w]}]$ గానీ లేదు.}}{}
    \tagitem{prvTrue}{\indcase{!A}{\ltrue}{$\mSat{M}{\indfrm}[w]$,
        $\mSat{M^*}{\indfrm}[{[w]}]$ రెండూ ఉన్నాయి.}}{}
    \item \indcase{!A}{p}{ఎడమ నుంచి కుడికి దిశ నేరుగా వస్తుంది.
      $\mSat{M}{\indfrm}[w]$ అయితే $w \in V(p)$ కావాలి.
      దానివల్ల $[w] \in V^*(p)$, అంటే
      $\mSat{M^*}{\indfrm}[{[w]}]$. తిరుగు దిశలో,
      $\mSat{M^*}{\indfrm}[{[w]}]$, అంటే
      $[w] \in V^*(p)$ అనుకుందాం. అప్పుడు ఏదో
      $v \in V(p)$కు $w \equiv v$. కనుక
      $\mSat{M}{p}[v]$ కూడా. $w \equiv v$ కనుక
      $w$, $v$లు $\Gamma$లోని ఒకే
      \tetoken{సూత్రాలను}{!!{formula}s} సత్యం చేస్తాయి.
      పరికల్పన ప్రకారం $p \in \Gamma$, అలాగే
      $\mSat{M}{p}[v]$ కనుక
      $\mSat{M}{\indfrm}[w]$.}

    \tagitem{prvNot}{\indcase{!A}{\lnot !B}
         {\iftag{probNot}{అభ్యాసం.}
         {$\mSat{M}{\indfrm}[w]$ అయితే, అప్పుడు మాత్రమే
          $\mSat/{M}{!B}[w]$. ఆగమన పరికల్పన ప్రకారం,
          $\mSat/{M}{!B}[w]$ అయితే, అప్పుడు మాత్రమే
          $\mSat/{M^*}{!B}[{[w]}]$. చివరగా,
          $\mSat/{M^*}{!B}[{[w]}]$ అయితే, అప్పుడు మాత్రమే
          $\mSat{M^*}{\indfrm}[{[w]}]$.}}}{}

    \tagitem{prvAnd}{\indcase{!A}{(!B \land !C)}
          {\iftag{probAnd}{అభ్యాసం.}
          {$\mSat{M}{\indfrm}[w]$ అయితే, అప్పుడు మాత్రమే
          $\mSat{M}{!B}[w]$, $\mSat{M}{!C}[w]$ రెండూ.
          ఆగమన పరికల్పన ప్రకారం, $\mSat{M}{!B}[w]$
          అయితే, అప్పుడు మాత్రమే $\mSat{M^*}{!B}[{[w]}]$;
          $\mSat{M}{!C}[w]$ అయితే, అప్పుడు మాత్రమే
          $\mSat{M^*}{!C}[{[w]}]$. అలాగే
          $\mSat{M^*}{\indfrm}[{[w]}]$ అయితే, అప్పుడు మాత్రమే
          $\mSat{M^*}{!B}[{[w]}]$,
          $\mSat{M^*}{!C}[{[w]}]$ రెండూ.}}}{}

    \tagitem{prvOr}{\indcase{!A}{(!B \lor !C)}
          {\iftag{probOr}{అభ్యాసం.}
          {$\mSat{M}{\indfrm}[w]$ అయితే, అప్పుడు మాత్రమే
          $\mSat{M}{!B}[w]$ లేదా $\mSat{M}{!C}[w]$.
          ఆగమన పరికల్పన ప్రకారం, $\mSat{M}{!B}[w]$
          అయితే, అప్పుడు మాత్రమే $\mSat{M^*}{!B}[{[w]}]$;
          $\mSat{M}{!C}[w]$ అయితే, అప్పుడు మాత్రమే
          $\mSat{M^*}{!C}[{[w]}]$. అలాగే
          $\mSat{M^*}{\indfrm}[{[w]}]$ అయితే, అప్పుడు మాత్రమే
          $\mSat{M^*}{!B}[{[w]}]$ లేదా
          $\mSat{M^*}{!C}[{[w]}]$.}}}{}

    \tagitem{prvIf}{\indcase{!A}{(!B \lif !C)}
          {\iftag{probIf}{అభ్యాసం.}
          {$\mSat{M}{\indfrm}[w]$ అయితే, అప్పుడు మాత్రమే
          $\mSat/{M}{!B}[w]$ లేదా $\mSat{M}{!C}[w]$.
          ఆగమన పరికల్పన ప్రకారం, $\mSat{M}{!B}[w]$
          అయితే, అప్పుడు మాత్రమే $\mSat{M^*}{!B}[{[w]}]$;
          $\mSat{M}{!C}[w]$ అయితే, అప్పుడు మాత్రమే
          $\mSat{M^*}{!C}[{[w]}]$. అలాగే
          $\mSat{M^*}{\indfrm}[{[w]}]$ అయితే, అప్పుడు మాత్రమే
          $\mSat/{M^*}{!B}[{[w]}]$ లేదా
          $\mSat{M^*}{!C}[{[w]}]$.}}}{}

    \tagitem{prvIff}{\indcase{!A}{(!B \liff !C)}
          {\iftag{probIff}{అభ్యాసం.}
          {$\mSat{M}{\indfrm}[w]$ అయితే, అప్పుడు మాత్రమే
          $\mSat{M}{!B}[w]$ మరియు $\mSat{M}{!C}[w]$,
          లేదా $\mSat/{M}{!B}[w]$ మరియు
          $\mSat/{M}{!C}[w]$. ఆగమన పరికల్పన ప్రకారం,
          $\mSat{M}{!B}[w]$ అయితే, అప్పుడు మాత్రమే
          $\mSat{M^*}{!B}[{[w]}]$;
          $\mSat{M}{!C}[w]$ అయితే, అప్పుడు మాత్రమే
          $\mSat{M^*}{!C}[{[w]}]$. అలాగే
          $\mSat{M^*}{\indfrm}[{[w]}]$ అయితే, అప్పుడు మాత్రమే
          $\mSat{M^*}{!B}[{[w]}]$ మరియు
          $\mSat{M^*}{!C}[{[w]}]$, లేదా
          $\mSat/{M^*}{!B}[{[w]}]$ మరియు
          $\mSat/{M^*}{!C}[{[w]}]$.}}}{}

    \tagitem{prvBox}{\indcase{!A}{\Box !B}
          {\iftag{probBox}{అభ్యాసం.}
          {$\mSat{M}{\indfrm}[w]$ అనుకుందాం.
          $\mSat{M^*}{\indfrm}[{[w]}]$ అని చూపడానికి,
          $R^*[w][v]$ అయ్యే $v$ను తీసుకుందాం.
          \olref{defn:filtration}\olref{defn:filtration-R2} వల్ల
          $\mSat{M}{!B}[v]$; ఆగమన పరికల్పన వల్ల
          $\mSat{M^*}{!B}[{[v]}]$. $v$ ఏదైనా కావచ్చు కనుక
          $\mSat{M^*}{\indfrm}[{[w]}]$.

          తిరుగు దిశలో, $\mSat{M^*}{\indfrm}[{[w]}]$ అనుకుందాం.
          $Rwv$ అయ్యే ఏదైనా $v$ను తీసుకుందాం.
          \olref{defn:filtration}\olref{defn:filtration-R1} వల్ల
          $R^*[w][v]$; అందువల్ల
          $\mSat{M^*}{!B}[{[v]}]$. ఆగమన పరికల్పన వల్ల
          $\mSat{M}{!B}[v]$; $v$ ఏదైనా కావచ్చు కనుక
          $\mSat{M}{\indfrm}[w]$.}}}{}

    \tagitem{prvDiamond}{\indcase{!A}{\Diamond !B}
          {\iftag{probDiamond}{అభ్యాసం.}
          {$\mSat{M}{\indfrm}[w]$ అనుకుందాం. అప్పుడు
          ఏదో $v \in W$కు $Rwv$, $\mSat{M}{!B}[v]$.
          ఆగమన పరికల్పన వల్ల $\mSat{M^*}{!B}[{[v]}]$;
          \olref{defn:filtration}\olref{defn:filtration-R1} వల్ల
          $R^*[w][v]$. కనుక $\mSat{M^*}{\indfrm}[{[w]}]$.

          ఇప్పుడు $\mSat{M^*}{\indfrm}[{[w]}]$ అనుకుందాం.
          అప్పుడు $R^*[w][v]$ అయ్యే ఏదో $[v]
          \in W^*$కు $\mSat{M^*}{!B}[{[v]}]$.
          ఆగమన పరికల్పన వల్ల $\mSat{M}{!B}[v]$.
          \olref{defn:filtration}\olref{defn:filtration-R3} వల్ల
          $\mSat{M}{\indfrm}[w]$.}}}{}
  \end{enumerate}
\end{proof}

\begin{prob}
  \olref[nml][fil][fil]{thm:filtrations} నిరూపణను పూర్తిచేయండి.
\end{prob}

ఒక లోకంలో సత్యానికి వర్తించే ఈ ఫలితం
నమూనాలో సత్యానికీ నమూనాల వర్గంలో చెల్లుబాటుకూ
వర్తిస్తుంది.

\begin{cor}
  $\Gamma$ ఉపసూత్రాల పరంగా సంవృతమై ఉండనివ్వండి.
  అప్పుడు:
  \begin{enumerate}
  \item $\mModel{M^*}$ అనేది $\Gamma$ ద్వారా
    $\mModel{M}$కు చేసే వడపోత అయితే, ఏ
    $!A \in \Gamma$కైనా $\mSat{M}{!A}$ కావడం
    $\mSat{M^*}{!A}$ కావడానికి తుల్యం.
  \item $\mClass{C}$ నమూనాల వర్గం,
    $\Gamma(\mClass{C})$ అనేది $\mClass{C}$లోని
    నమూనాలకు చేసే $\Gamma$-వడపోతల వర్గం అయితే,
    ఏ \tetoken{సూత్రం}{!!{formula}}~$!A \in \Gamma$ అయినా
    $\mClass{C}$లో చెల్లుబాటు కావడం
    $\Gamma(\mClass{C})$లో చెల్లుబాటు కావడానికి తుల్యం.
  \end{enumerate}
\end{cor}



\end{document}
